\documentclass{article}

\usepackage[italian]{babel}

\usepackage[a4paper,top=2cm,bottom=2cm,left=3cm,right=3cm,marginparwidth=1.75cm]{geometry}

\title{Buckshot Roulette}
\author{Alpi Alessandro, Casadei Veronica, Conca Geremia, Fiumana Giulia}

\begin{document}
\maketitle

\tableofcontents

\section{Analisi}
\subsection{Descrizione e requisiti}
Il software, commissionato dal corso "Programmazione ad Oggetti" del corso "Ingegneria e scienze informatiche", mira alla creazione di un'intelligenza artificiale denominata The Dealer (Il Mazziere). Per intelligenza artificiale si intende un software in grado di assumere decisioni complesse in maniera semi-autonoma sugli argomenti di sua competenza, a partire dai vincoli e dagli obiettivi datigli dall'utente. In questo caso, l'intelligenza artificiale dovrà interagire con l'utente e sfidarlo ad una serie di partite di Buckshot Roulette, una variante della roulette russa. Durante il gioco saranno usati oggetti e pulsanti. La partita termina alla morte del dealer o dell'utente, con la possibilità di rigiocare a scelta del giocatore. In caso di vincita dell'utente, questo otterrà una somma di denaro.


\subsection{Modello del Dominio}


\section{Design}
\subsection{Architettura}
\subsection{Design dettagliato}

\section{Sviluppo}
\subsection{Testing automatizzato}
\subsection{Note di sviluppo}
\subsubsection{Esempio}

\section{Commenti finali}
\subsection{Autovalutazione e lavori futuri}
\subsection{Difficoltà incontrate e commenti per i docenti}

\appendix
\section{Guida utente}
\section{Esercitazioni di laboratorio}



\end{document}
